% ============================================================
% chapters/00-abstract.tex -- תקציר מנהלים ומילות מפתח
% Executive summary: 250-350 words, one continuous block (booklet, sec. 2.2).
% Written early as a draft; REVISE on Monday after chapters 6-8 are final.
% Content only. The template decides where and how it is printed.
% Template must provide: environment `executivesummary`, command \keywords{...}
% ============================================================

\begin{executivesummary}
בשנת 410 בזזו הוויזיגותים את רומא, ומאז ועד היום חוזרת השאלה אם קבלת הנצרות החלישה את הקיסרות הרומית. גיבון השיב בחיוב, ואילו רוב החוקרים בני זמננו מדגישים לחץ חיצוני, משבר כלכלי, אקלים ומגפות. למרות מאות שנות מחקר, הוויכוח לא הוכרע: רוב המחקרים בוחנים מקרה אחד בכל פעם, נשענים על דוגמאות נבחרות, ואינם מגדירים את המושגים שלהם באופן מדיד. דוח מכין זה מציע מחקר שבוחן את השאלה בדרך חדשה, באמצעות השוואה בין שלוש מדינות נוצריות שהגיעו לתוצאות שונות לגמרי: רומא המערבית, שנפלה ב־476; רומא המזרחית, שאימצה את אותה דת ושרדה עוד כאלף שנה; והממלכה הפרנקית, שהייתה נוצרית כמעט מהקמתה והתחזקה. שאלת המחקר היא אם הנצרות פעלה ככוח שמחליש או שמחזק מדינות בין המאה הרביעית למאה התשיעית, ומה קבע לאיזה כיוון היא פעלה. ההשערה המרכזית היא שהדת עצמה לא הכריעה, אלא מידת ההשתלבות בין הכנסייה לבין מוסדות המדינה. לאחר השוואה בין שלוש חלופות שיטה במטריצת החלטה משוקללת, נבחרה שיטה משולבת: מאגר נתונים השוואתי במקטעים של 25 שנה, שבו מדדי ההשתלבות מופקים מאלפי טקסטים עתיקים בעזרת מודלי שפה וצוות של סוכני \en{AI}, ונבדקים בידי מומחה לתקופה. המחקר נשען על כלים קיימים בגישה פתוחה – ספריות הטקסטים \en{Perseus} ו־\en{dMGH} ומאגר \en{Seshat} – ומוסיף להם שכבה שחסרה בהם: מדידה שיטתית של היחסים בין הכנסייה למדינה. מפרט הדרישות קובע, בין היתר, שסיווגי המודל יתאימו לפחות ב־85\% לסיווג של מומחה, ושכל תוצאה תהיה ניתנת לשחזור מהקוד ומהנתונים. ניתוח רגישות הראה שהבחירה יציבה גם כשהמשקלים משתנים. המחקר מתוכנן לשנה, בחמש חבילות עבודה: בניית הקורפוס, בניית המאגר, ניתוח הטקסטים ותיקופו, ניתוח סטטיסטי של ההשערות, ופרסום פתוח של הקוד, הנתונים והממצאים. הסיכונים העיקריים, ובראשם הטיית המקורות שנכתבו ברובם בידי אנשי כנסייה וטעויות של מודלי השפה, מטופלים בתוכנית הפחתה מוגדרת. תוצרי המחקר יהיו מאגר וקורפוס מתויג בגישה פתוחה, שיאפשרו לראשונה לבחון את הוויכוח בן 1,600 השנים כשאלה אמפירית.
\end{executivesummary}

\keywords{נפילת רומא; התנצרות הקיסרות הרומית; הממלכה הפרנקית; יחסי דת ומדינה; יציבות מדינית; היסטוריה ממוחשבת; ניתוח טקסט בעזרת \en{AI}}
